% ------------------------------------------------------------
% LaTeX Template für die DHBW zum Schnellstart!
% Original: https://github.wdf.sap.corp/vtgermany/LaTeX-Template-DHBW
% ------------------------------------------------------------
% ---- Präambel mit Angaben zum Dokument
\documentclass[
	fontsize=12pt,           % Leitlinien sprechen von Schriftgröße 12.
	paper=A4,
	twoside=false,
	listof=totoc,            % Tabellen- und Abbildungsverzeichnis ins Inhaltsverzeichnis
	bibliography=totoc,      % Literaturverzeichnis ins Inhaltsverzeichnis aufnehmen
	titlepage,               % Titlepage-Umgebung anstatt \maketitle
	headsepline,             % horizontale Linie unter Kolumnentitel
	abstracton,              % Überschrift einschalten, Abstract muss in {abstract}-Umgebung stehen
]{scrreprt}                  % Verwendung von KOMA-Report
\usepackage[utf8]{inputenc}  % UTF8 Encoding einschalten
\usepackage[english]{babel}  % Neue deutsche Rechtschreibung
\usepackage[T1]{fontenc}     % Ausgabe von westeuropäischen Zeichen (auch Umlaute)
\usepackage{lmodern}         % Nicht-gerasterte Schriftarten (bei MikTeX erforderlich)
\usepackage{graphicx}        % Einbinden von Grafiken erlauben
\usepackage{wrapfig}         % Grafiken fließend im Text
\usepackage{setspace}        % Zeilenabstand \singlespacing, \onehalfspaceing, \doublespacing
\usepackage[
	%showframe,                % Ränder anzeigen lassen
	left=2.7cm, right=2.5cm,
	top=2.5cm,  bottom=2.5cm,
	includeheadfoot
]{geometry}                      % Seitenlayout einstellen
\usepackage{scrlayer-scrpage}   
\usepackage{pgfplots}% Gestaltung von Fuß- und Kopfzeilen
\usepackage{acronym}   
\usepackage{mathrsfs}          % Abkürzungen, Abkürzungsverzeichnis
\usepackage{titletoc}            % Anpassungen am Inhaltsverzeichnis
\contentsmargin{0.725cm}         % Abstand im Inhaltsverzeichnis zw. Punkt und Seitenzahl
\usepackage[hyphens]{url}
\usepackage[                     % Klickbare Links (enth. auch "nameref", "url" Package)
  hidelinks,                     % Blende die "URL Boxen" aus.
  breaklinks=true                % Breche zu lange URLs am Zeilenende um
]{hyperref}
\urlstyle{same}                  % Aktuelle Schrift auch für URLs
% Anpassung von autoref für Gleichungen (ergänzt runde Klammern) und Algorithm.
% Anstatt "Listing" kann auch z.B. "Code-Ausschnitt" verwendet werden. Dies sollte
% jedoch synchron gehalten werden mit \lstlistingname (siehe weiter unten).
\addto\extrasenglish{%
	\def\equationautorefname~#1\null{Equation~(#1)\null}
	\def\lstnumberautorefname{Line}
	\def\lstlistingautorefname{Listing}
	\def\algorithmautorefname{Algorithm}
	% Damit einheitlich "Abschnitt 1.2[.3]" verwendet wird und nicht "Unterabschnitt 1.2.3"
	% \def\subsectionautorefname{Abschnitt}
}

% ---- Abstand verkleinern von der Überschrift 
\renewcommand*{\chapterheadstartvskip}{\vspace*{.5\baselineskip}}

% Hierdurch werden Schusterjungen und Hurenkinder vermieden, d.h. einzelne Wörter
% auf der nächsten Seite oder in einer einzigen Zeile.
% LaTeX kann diese dennoch erzeugen, falls das Layout ansonsten nicht umsetzbar ist.
% Diese Werte sind aber gute Startwerte.
\widowpenalty10000
\clubpenalty10000
% ---- Für das Quellenverzeichnis
\usepackage[
	backend = biber,                % Verweis auf biber
	language = auto,
	style = numeric,                % Nummerierung der Quellen mit Zahlen
	sorting = none,                 % none = Sortierung nach der Erscheinung im Dokument
	sortcites = true,               % Sortiert die Quellen innerhalb eines cite-Befehls
	block = space,                  % Extra Leerzeichen zwischen Blocks
	hyperref = true,                % Links sind klickbar auch in der Quelle
	%backref = true,                % Referenz, auf den Text an die zitierte Stelle
	bibencoding = auto,
	giveninits = true,              % Vornamen werden abgekürzt
	doi=false,                      % DOI nicht anzeigen
	isbn=false,                     % ISBN nicht anzeigen
    alldates=short                  % Datum immer als DD.MM.YYYY anzeigen
]{biblatex}
\addbibresource{content/literatur.bib}
\setcounter{biburlnumpenalty}{3000}     % Umbruchgrenze für Zahlen
\setcounter{biburlucpenalty}{6000}      % Umbruchgrenze für Großbuchstaben
\setcounter{biburllcpenalty}{9000}      % Umbruchgrenze für Kleinbuchstaben
\DeclareNameAlias{default}{family-given}  % Nachname vor dem Vornamen
\AtBeginBibliography{\renewcommand{\multinamedelim}{\addslash\space
}\renewcommand{\finalnamedelim}{\multinamedelim}}  % Schrägstrich zwischen den Autorennamen
\DefineBibliographyStrings{german}{
  urlseen = {Einsichtnahme:},                      % Ändern des Titels von "besucht am"
}
\usepackage[babel,german=quotes]{csquotes}         % Deutsche Anführungszeichen + Zitate


% ---- Für Mathevorlage
\usepackage{amsmath}    % Erweiterung vom Mathe-Satz
\usepackage{amssymb}    % Lädt amsfonts und weitere Symbole
\usepackage{MnSymbol}   % Für Symbole, die in amssymb nicht enthalten sind.


% ---- Für source-codevorlage
\usepackage{scrhack}                    % Hack zur Verw. von listings in KOMA-Script
\usepackage{listings}                   % Darstellung von source-code
\usepackage{xcolor}                     % Einfache Verwendung von Farben

% -- Eigene Farben für den source-code
\definecolor{JavaLila}{rgb}{0.4,0.1,0.4}
\definecolor{JavaGruen}{rgb}{0.3,0.5,0.4}
\definecolor{JavaBlau}{rgb}{0.0,0.0,1.0}
\definecolor{ABAPKeywordsBlue}{HTML}{6000ff}
\definecolor{ABAPCommentGrey}{HTML}{808080}
\definecolor{ABAPStringGreen}{HTML}{4da619}
\definecolor{PyKeywordsBlue}{HTML}{0000AC}
\definecolor{PyCommentGrey}{HTML}{808080}
\definecolor{PyStringGreen}{HTML}{008080}

% -- Default Listing-Styles

\lstset{
	% Das Paket "listings" kann kein UTF-8. Deswegen werden hier 
	% die häufigsten Zeichen definiert (ä,ö,ü,...)
	literate=%
		{á}{{\'a}}1 {é}{{\'e}}1 {í}{{\'i}}1 {ó}{{\'o}}1 {ú}{{\'u}}1
		{Á}{{\'A}}1 {É}{{\'E}}1 {Í}{{\'I}}1 {Ó}{{\'O}}1 {Ú}{{\'U}}1
		{à}{{\`a}}1 {è}{{\`e}}1 {ì}{{\`i}}1 {ò}{{\`o}}1 {ù}{{\`u}}1
		{À}{{\`A}}1 {È}{{\'E}}1 {Ì}{{\`I}}1 {Ò}{{\`O}}1 {Ù}{{\`U}}1
		{ä}{{\"a}}1 {ë}{{\"e}}1 {ï}{{\"i}}1 {ö}{{\"o}}1 {ü}{{\"u}}1
		{Ä}{{\"A}}1 {Ë}{{\"E}}1 {Ï}{{\"I}}1 {Ö}{{\"O}}1 {Ü}{{\"U}}1
		{â}{{\^a}}1 {ê}{{\^e}}1 {î}{{\^i}}1 {ô}{{\^o}}1 {û}{{\^u}}1
		{Â}{{\^A}}1 {Ê}{{\^E}}1 {Î}{{\^I}}1 {Ô}{{\^O}}1 {Û}{{\^U}}1
		{œ}{{\oe}}1 {Œ}{{\OE}}1 {æ}{{\ae}}1 {Æ}{{\AE}}1 {ß}{{\ss}}1
		{ű}{{\H{u}}}1 {Ű}{{\H{U}}}1 {ő}{{\H{o}}}1 {Ő}{{\H{O}}}1
		{ç}{{\c c}}1 {Ç}{{\c C}}1 {ø}{{\o}}1 {å}{{\r a}}1 {Å}{{\r A}}1
		{€}{{\euro}}1 {£}{{\pounds}}1 {«}{{\guillemotleft}}1
		{»}{{\guillemotright}}1 {ñ}{{\~n}}1 {Ñ}{{\~N}}1 {¿}{{?`}}1,
	breaklines=true,        % Breche lange Zeilen um 
	breakatwhitespace=true, % Wenn möglich, bei Leerzeichen umbrechen
	% Symbol für Zeilenumbruch einfügen
	prebreak=\raisebox{0ex}[0ex][0ex]{\ensuremath{\rhookswarrow}},
	postbreak=\raisebox{0ex}[0ex][0ex]{\ensuremath{\rcurvearrowse\space}},
	tabsize=4,                                 % Setze die Breite eines Tabs
	basicstyle=\ttfamily\small,                % Grundsätzlicher Schriftstyle
	columns=fixed,                             % Besseres Schriftbild
	numbers=left,                              % Nummerierung der Zeilen
	%frame=single,                             % Umrandung des Codes
	showstringspaces=false,                    % Keine Leerzeichen hervorheben
	keywordstyle=\color{blue},
	ndkeywordstyle=\bfseries\color{darkgray},
	identifierstyle=\color{black},
	commentstyle=\itshape\color{JavaGruen},   % Kommentare in eigener Farbe
	stringstyle=\color{JavaBlau},             % Strings in eigener Farbe,
	captionpos=b,                             % Bild*unter*schrift
	xleftmargin=5.0ex
}

% ---- Eigener JAVA-Style für den source-code
\renewcommand{\ttdefault}{pcr}               % Schriftart, welche auch fett beinhaltet
\lstdefinestyle{EigenerJavaStyle}{
	language=Java,                             % Syntax Highlighting für Java
	%frame=single,                             % Umrandung des Codes
	keywordstyle=\bfseries\color{JavaLila},    % Keywords in eigener Farbe und fett
	commentstyle=\itshape\color{JavaGruen},    % Kommentare in eigener Farbe und italic
	stringstyle=\color{JavaBlau}               % Strings in eigener Farbe
}

% ---- Eigener ABAP-Style für den source-code
\renewcommand{\ttdefault}{pcr}
\lstdefinestyle{EigenerABAPStyle}{
	language=[R/3 6.10]ABAP,
	morestring=[b]\|,                          % Für Pipe-Strings
	morestring=[b]\`,                          % für Backtick-Strings
	keywordstyle=\bfseries\color{ABAPKeywordsBlue},
	commentstyle=\itshape\color{ABAPCommentGrey},
	stringstyle=\color{ABAPStringGreen},
	tabsize=2
}

% ---- Eigener Python-Style für den source-code
\renewcommand{\ttdefault}{pcr}
\lstdefinestyle{EigenerPythonStyle}{
	language=Python,
	columns=flexible,
	keywordstyle=\bfseries\color{PyKeywordsBlue},
	commentstyle=\itshape\color{PyCommentGrey},
	stringstyle=\color{PyStringGreen}
}  % Weitere Details sind ausgelagert

\usepackage{algorithm}                  % Für Algorithmen-Umgebung (ähnlich wie lstlistings Umgebung)
\usepackage{algpseudocode}              % Für Pseudocode. Füge "[noend]" hinzu, wenn du kein "endif",
                                        % etc. haben willst.

\makeatletter                           % Sorgt dafür, dass man @ in Namen verwenden kann.
                                        % Ansonsten gibt es in der nächsten Zeile einen Compilefehler.
\renewcommand{\ALG@name}{Algorithm}   % Umbenennen von "Algorithm" im Header der Listings.
\makeatother                            % Zeichen wieder zurücksetzen
\renewcommand{\lstlistingname}{Listing} % Erlaubt das Umbenennen von "Listing" in anderen Titel.

% ---- Tabellen
\usepackage{booktabs}  % Für schönere Tabellen. Enthält neue Befehle wie \midrule
\usepackage{multirow}  % Mehrzeilige Tabellen
\usepackage{siunitx}   % Für SI Einheiten und das Ausrichten Nachkommastellen
\sisetup{locale=DE, range-phrase={~bis~}, output-decimal-marker={,}} % Damit ein Komma und kein Punkt verwendet wird.

% ---- Für Definitionsboxen in der Einleitung
\usepackage{amsthm}                     % Liefert die Grundlagen für Theoreme
\usepackage[framemethod=tikz]{mdframed} % Boxen für die Umrandung
% ---- Definition für Highlight Boxen

% ---- Grundsätzliche Definition zum Style
\newtheoremstyle{defi}
  {\topsep}         % Abstand oben
  {\topsep}         % Abstand unten
  {\normalfont}     % Schrift des Bodys
  {0pt}             % Einschub der ersten Zeile
  {\bfseries}       % Darstellung von der Schrift in der Überschrift
  {:}               % Trennzeichen zwischen Überschrift und Body
  {.5em}            % Abstand nach dem Trennzeichen zum Body Text
  {\thmname{#3}}    % Name in eckigen Klammern
\theoremstyle{defi}

% ------ Definition zum Strich vor eines Texts
\newmdtheoremenv[
  hidealllines = true,       % Rahmen komplett ausblenden
  leftline = true,           % Linie links einschalten
  innertopmargin = 0pt,      % Abstand oben
  innerbottommargin = 4pt,   % Abstand unten
  innerrightmargin = 0pt,    % Abstand rechts
  linewidth = 3pt,           % Linienbreite
  linecolor = gray!40,       % Linienfarbe
]{defStrich}{Definition}     % Name der des formats "defStrich"

% ------ Definition zum Eck-Kasten um einen Text
\newmdtheoremenv[
  hidealllines = true,
  innertopmargin = 6pt,
  linecolor = gray!40,
  singleextra={              % Eck-Markierungen für die Definition
    \draw[line width=3pt,gray!50,line cap=rect] (O|-P) -- +(1cm,0pt);
    \draw[line width=3pt,gray!50,line cap=rect] (O|-P) -- +(0pt,-1cm);
    \draw[line width=3pt,gray!50,line cap=rect] (O-|P) -- +(-1cm,0pt);
    \draw[line width=3pt,gray!50,line cap=rect] (O-|P) -- +(0pt,1cm);
  }
]{defEckKasten}{Definition}  % Name der des formats "defEckKasten"  % Weitere Details sind ausgelagert

% ---- Für Todo Notes
\usepackage{todonotes}


% ---- Elektronische Version oder Gedruckte Version?
% ---- Unterschied: Die elektronische Version enthält keinen Platzhalter für die Unterschrift
\usepackage{ifthen}
\newboolean{e-Abgabe}
\setboolean{e-Abgabe}{false}    % false=gedruckte Fassung

% ---- Persönlichen Daten:
\newcommand{\titel}{Template \LaTeX\ Wiki von BAzubis für BAzubis}
\newcommand{\titelheader}{Titel welcher im Header auftaucht}
\newcommand{\arbeit}{Projektarbeit 1 (T2\_2000)}
\newcommand{\studiengang}{Angewandte Informatik}
\newcommand{\studienjahr}{2015}
\newcommand{\autor}{Vorname Nachname}
\newcommand{\autorReverse}{Nachname, Vorname}
\newcommand{\verfassungsort}{Karlsruhe}
\newcommand{\matrikelnr}{0000000}
\newcommand{\kurs}{TINF15B1}
\newcommand{\bearbeitungsmonat}{Januar 2018}
\newcommand{\abgabe}{01. Februar 2018}
\newcommand{\bearbeitungszeitraum}{01.10.2017 - 31.01.2018}
\newcommand{\firmaName}{SAP SE}
\newcommand{\firmaStrasse}{Dietmar-Hopp-Allee 16}
\newcommand{\firmaPlz}{69190 Walldorf, Deutschland}
\newcommand{\betreuerFirma}{B-Vorname B-Nachname}
\newcommand{\betreuerDhbw}{DH-Vorname DH-Nachname}

% ---- Metainformation für das PDF Dokument
\hypersetup{
	pdftitle    = {\titel},
	pdfsubject  = {\arbeit},
	pdfauthor   = {\autor},
	%pdfkeywords = {Keywords angeben},
	pdfcreator  = {LaTeX},
	%pdfproducer = {in der Regel pdfTeX}
}

% ---- Definition der Kopf- und Fußzeilen
\clearscrheadfoot                               % Löschen von LaTeX Standard
\automark[section]{chapter}                     % Füllen von section und chapter
\renewcommand*{\chaptermarkformat}{}            % Entfernt die Kapitelnummer
\renewcommand*{\sectionmarkformat}{}            % Entfernt die Sectionnummer
% Angaben [für "plain"]{für "scrheadings"}
\ihead[]{\titelheader}                          % Kopfzeile links
\chead[]{}                                      % Kopfzeile mitte
\ohead[]{\rightmark}                            % Kopfzeile rechts
\ifoot[]{}                                      % Fußzeile links
\cfoot*{\sffamily\pagemark}                     % Fußzeile mitte
\ofoot[]{}                                      % Fußzeile rechts
\KOMAoptions{
   headsepline = 0.2pt,                         % Liniendicke Kopfzeile
   footsepline = false                          % Liniendicke Fußzeile
}

% ---- Hilfreiches
\newcommand{\zB}{z.\,B. }   % "z.B." mit kleinem Leeraum dazwischen (ohne wäre nicht korrekt)
\newcommand{\dash}{d.\,h. }

\newcommand{\code}[1]{\texttt{#1}} % Ist einfacher zu schreiben als ständig \texttt und erlaubt
                                   % Änderungen im Nachhinein, wenn man z.B. Inline-Code anders stylen möchte.

% ---- Silbentrennung (falls LaTeX defaults falsch / nicht gewünscht sind)
\hyphenation{HANA}         % anstatt HA-NA
\hyphenation{Graph-Script} % anstatt GraphS-cript

% ---- Beginn des Dokuments
\begin{document}
\setlength{\parindent}{0pt}              % Keine Paragraphen Einrückung.
                                         % Dafür haben wir den Abstand zwischen den Paragraphen.
\setcounter{secnumdepth}{2}              % Nummerierungstiefe fürs Inhaltsverzeichnis
\setcounter{tocdepth}{1}                 % Tiefe des Inhaltsverzeichnisses. Ggf. so anpassen,
                                         % dass das Verzeichnis auf eine Seite passt.
\sffamily                                % Serifenlose Schrift verwenden.

% ---- Vorspann
% ------ Titelseite
\singlespacing
\thispagestyle{empty}
\begin{titlepage}
\enlargethispage{4cm}

\begin{figure}           % Logo vom Ausbildungsbetrieb und der DHBW
	% \vspace*{-5mm} % Sollte dein Titel zu lang werden, kannst du mit diesem "Hack" 
	%                  den Inhalt der Seite nach oben schieben.
	\begin{minipage}{0.49\textwidth}
		\flushleft
		\includegraphics[height=2.5cm]{images/Logos/Logo_DHBW.pdf} 
	\end{minipage}
\end{figure} 
\vspace*{0.1cm}

\begin{center}
	\huge{\textbf{\titel}}\\[1.5cm]
	\Large{\textbf{\arbeit}}\\[0.5cm]
	\normalsize{im Rahmen der Prüfung zum\\[1ex] \textbf{Bachelor of Science (B.Sc.)}}\\[0.5cm]
	\Large{des Studienganges \studiengang}\\[1ex]
	\normalsize{an der Dualen Hochschule Baden-Württemberg Karlsruhe}\\[1cm]
	\normalsize{von}\\[1ex] \Large{\textbf{\autor}} \\[1cm]
	\normalsize{\bearbeitungsmonat}\\[1ex] \Large{\textbf{-Sperrvermerk-}}\\[0.5cm]
\end{center}

\begin{center}
	\vfill
	\begin{tabular}{ll}
		Abgabedatum:                     & \abgabe \\[0.2cm]
		Bearbeitungszeitraum:            & \bearbeitungszeitraum \\[0.2cm]
		Matrikelnummer, Kurs:            & \matrikelnr , \kurs \\[0.2cm]
		Ausbildungsfirma:                & \firmaName \\
		                                 & \firmaStrasse \\
		                                 & \firmaPlz \\[0.2cm]
		Betreuer der Ausbildungsfirma:   & \betreuerFirma \\[0.2cm]
		Gutachter der Dualen Hochschule: & \betreuerDhbw \\[2cm]
	\end{tabular} 
\end{center}
\end{titlepage}
  % Titelseite
\newcounter{savepage}
\pagenumbering{Roman}                    % Römische Seitenzahlen
\onehalfspacing

% ------ Erklärung, Sperrvermerk, Abstact
\chapter*{Eidesstattliche Erklärung}
Ich versichere hiermit, dass ich meine \arbeit{} mit dem Thema:
\begin{quote}
	\textit{\titel}
\end{quote} 
gemäß § 5 der \enquote{Studien- und Prüfungsordnung DHBW Technik} vom 29. September 2017 selbstständig verfasst und keine anderen als die angegebenen Quellen und Hilfsmittel benutzt habe. Die Arbeit wurde bisher keiner anderen Prüfungsbehörde vorgelegt und auch nicht veröffentlicht.

\vspace{0.25cm}

Ich versichere zudem, dass die eingereichte elektronische Fassung mit der gedruckten Fassung übereinstimmt.

\vspace{1cm}

\verfassungsort, den \today \\[0.5cm]
\ifthenelse{\boolean{e-Abgabe}}
	{\underline{Gez. \autor}}
	{\makebox[6cm]{\hrulefill}}\\ 
\autorReverse

\chapter*{Sperrvermerk}
Die nachfolgende Arbeit enthält vertrauliche Daten der:
\begin{quote}
	\firmaName \\
	\firmaStrasse \\
	\firmaPlz
\end{quote}

\vspace{0.5cm}

Sie darf als Leistungsnachweis des Studienganges \studiengang{} \studienjahr{} an der DHBW Karlsruhe verwendet und nur zu Prüfungszwecken zugänglich gemacht werden. Über den Inhalt ist Stillschweigen zu bewahren. Veröffentlichungen oder Vervielfältigungen der \arbeit{} - auch auszugsweise - sind ohne ausdrückliche Genehmigung der SAP SE nicht gestattet.

\vspace{0.5cm}

SAP und die SAP Logos sind eingetragene Warenzeichen der SAP SE.
Die Wiedergabe von Gebrauchsnamen, Handelsnamen, Warenbezeichnungen usw. in dieser Arbeit berechtigt auch ohne besondere Kennzeichnung nicht zu der Annahme, dass solche Namen im Sinne der Warenzeichen- und Markenschutz-Gesetzgebung als frei zu betrachten wären und daher von jedem benutzt werden dürfen.

\renewcommand{\abstractname}{Abstract} % Veränderter Name für das Abstract
\begin{abstract}
\begin{addmargin}[1.5cm]{1.5cm}        % Erhöhte Ränder, für Abstract Look
\thispagestyle{plain}                  % Seitenzahl auf der Abstract Seite

\begin{center}
\small\textit{- English -}             % Angabe der Sprache für das Abstract
\end{center}

\vspace{0.25cm}

This is the starting point of the Abstract. For the final bachelor thesis, there must be an abstract included in your document. So, start now writing it in German and English. The abstract is a short summary with around 200 to 250 words.

\vspace{0.25cm}

Try to include in this abstract the main question of your work, the methods you used or the main results of your work.


\end{addmargin}
\end{abstract}
\renewcommand{\abstractname}{Abstract} % Veränderter Name für das Abstract
\begin{abstract}
\begin{addmargin}[1.5cm]{1.5cm}        % Erhöhte Ränder, für Abstract Look
\thispagestyle{plain}                  % Seitenzahl auf der Abstract Seite
This paper presents and compares four different model order reduction methods.
These are proper orthogonal decomposition, balanced truncation, modal truncation, and Hankel norm approximation.
These methods are applied to the 1D heat equation.
The heat equation is modified so that a heating process can be represented.
This equation is discretized using the finite element method to obtain a state space representation of the system.
The model order reduction methods are compared with respect to the \(L2\)- and \(H_{\infty}\) error and the runtime.

\vspace{0.25cm}



\end{addmargin}
\end{abstract}

% ------ Inhaltsverzeichnis
\singlespacing
\tableofcontents

% ------ Verzeichnisse
\renewcommand*{\chapterpagestyle}{plain}
\pagestyle{plain}
\chapter*{List of abbreviations}
\addcontentsline{toc}{chapter}{Abkürzungsverzeichnis} % Hinzufügen zum Inhaltsverzeichnis 

\begin{acronym}[Inverse Discrete Fourier Transform] % längstes Kürzel wird verw. für den Abstand zw. Kürzel u. Text
	\acro{DE}{Differential Equation}
	\acro{DFT}{Discrete Fourier Transform}
	\acro{FEM}{Finite Element Method}
	\acro{FT}{Fourier Transform}	
	\acro{FFT}{Fast Fourier Transform}
	\acro{IDFT}{Inverse Discrete Fourier Transform}
	\acro{IFFT}{Inverse Fast Fourier Transform}
	\acro{LS}{Least Squares}
	\acro{ODE}{Ordinary Differential Equation}
	\acro{PDE}{Partial Differential Equation}
\end{acronym}

\listoffigures                          % Erzeugen des Abbildungsverzeichnisses 
\listoftables                           % Erzeugen des Tabellenverzeichnisses
\renewcommand{\lstlistlistingname}{source-codeverzeichnis}
\lstlistoflistings                      % Erzeugen des Listenverzeichnisses
\setcounter{savepage}{\value{page}}


% ---- Inhalt der Arbeit
\cleardoublepage
\pagenumbering{arabic}                  % Arabische Seitenzahlen für den Hauptteil
\chapter{Fundamentals}
\section{Heat Equation} \label{HE}
The conduction of heat  within a medium can be described using the following partial differential equation (PDE):
\begin{gather}
\frac{\partial u}{\partial t} = \alpha \nabla^{2} u \label{heat}
\end{gather}
With \(u\) being a function of space and time and \(\alpha\) being a positive constant.
For this paper \(u\) will be defined in terms of one spacial dimension:
\begin{gather}
u := u(x, t) \\
\frac{\partial u}{\partial t} = \alpha \frac{\partial^{2} u}{\partial x^{2}} \\
x \in \chi \subset \mathbb{R} \quad t \in \tau \subset \mathbb{R} \\
x_{0} \leq x \leq x_{n} \quad t_{0} \leq t \leq t_{n} 
\end{gather}
\cite{Gustafsson2011}

In order to not only model the conduction of heat within a medium but also a heating process, a new function \(h: \chi \times \tau \rightarrow  \mathbb{R}\) is introduced:
\begin{gather}
\frac{\partial u}{\partial t} = \alpha \frac{\partial^{2} u}{\partial x^{2}} + h(x,t) \label{eq-1d-h}
\end{gather}
For this paper it is assumed that the initial condition is known:
\begin{gather}
f: \chi \rightarrow \mathbb{R} \\
u(x, t_{0}) = f(x) 
\end{gather} 
\subsection{Solving the Heat Equation using Fourier transform}
Applying the Fourier transform (FT) w.r.t \(x\) to \ref{eq-1d-h} yields the inhomogeneous ordinary differential equation (ODE):

\begin{gather}
\hat{u} = \mathfrak{F}(u) \quad \hat{h} = \mathfrak{F}(h) \\
\frac{d}{dt} \hat{u} = -\alpha\omega^{2}\hat{u} + \hat{h} \label{eq-1d-h-ft}
\end{gather}

A solution to \ref{eq-1d-h-ft} is given by:
\begin{gather}
\hat{u} = \hat{u}_{0} + \hat{u}_{p}
\end{gather}
Where \(\hat{u}_{0}\) is the homogeneous solution and \(\hat{u}_{p}\) is the particular integral.
In order to solve this ODE for the particular integral \(\hat{h}\) has to be known. \cite{Papula2015}
The choice of \(h\) is, except to some restrictions, arbitrary.
Therefore an approximate solution to \ref{eq-1d-h-ft} \(\hat{u}_{a}\) is obtained by the forward euler scheme:
\begin{gather}
\frac{d}{dt} \hat{u} \approx \frac{\Delta \hat{u}}{\Delta t} \\
\hat{u}_{t+1} = \hat{u}_{t} + \Delta t (-\alpha\omega^{2}\hat{u} + \hat{h}) \label{eq-1d-h-es} \\
\hat{u}_{a} = [\hat{u}_{t_{0}}, ..., \hat{u}_{t_{n}}]
\end{gather}
In order to apply the euler scheme successfully an initial condition   \(\hat{u}_{0}\) has to be known. 
This initial condition is obtained by applying the discrete Fourier transform (DFT) to an initial temperature distribution along \(x\):
\begin{gather}
\hat{u}_{0} = \mathfrak{F}(f(x)) 
\end{gather} 
\cite{Gustafsson2011b}

The forward euler scheme is used here because it is fairly easy to implement.
By applying the  inverse discrete Fourier transform (IDFT) to \(\hat{u}_{a}\) an approximate solution to \ref{eq-1d-h} can be obtained.
To decrease computing time, the fast Fourier transform (FFT) and inverse fast Fourier transform (IFFT) is used instead of the DFT and IDFT.

\section{Finite Element Method} \label{FEM}
The finite element method (FEM) is a method to approximate solutions for differential equations (DE) within a certain domain \(\Omega\).
This is done by discretizing the spatial domain.
Assume that a DE is given by:
\begin{gather}
m, n \in \mathbb{N} \quad \zeta \in \Omega \subset \mathbb{R} \quad m \geq 1 \\
\frac{\partial^{m} y}{\partial \zeta^{m}} -  g(y) = r(\zeta, t) \label{eq-pde-gen} 
\end{gather}
It is assumed that \(g\) is a linear function that can also contain partial derivatives of \(y\) w.r.t. time, \(y\) takes the value 0 at the boundary \(\Gamma\) and \(y(\zeta, 0) = f(\zeta)\).
An approximate solution to \(y\) is given by \(\mu\), which is expressed as a sum of basis functions contained in the set \(\phi\):
\begin{gather}
\mu(\zeta, t) = \sum_{j = 1}^{N} c_{j}(t)\phi_{j}(\zeta) \label{eq-aprox}
\end{gather}
The residual is defined as:
\begin{gather}
\mathfrak{r} = \frac{\partial^{m} \mu}{\partial \zeta^{m}} -  g(\mu) - r(\zeta, t) 
\end{gather}
Furthermore the residual is required to be orthogonal to all basis functions:
\begin{gather}
\langle \mathfrak{r}, \phi_{k} \rangle = 0 \quad \forall \phi_{k} \in \phi \label{eq-req-orth}
\end{gather}

Since the functions in \(\phi\) are known, it is only required to find the coefficients \(c_{j}(t)\) in \ref{eq-aprox}.
To find those coefficients \ref{eq-req-orth} needs to be expressed as follows:
\begin{gather}
\int_{\Omega} \frac{\partial^{m} \mu}{\partial \zeta^{m}} \phi_{k} \, d\zeta  -  \int_{\Omega} g(\mu) \phi_{k}        \, d\zeta = \int_{\Omega}  r(\zeta, t) \phi_{k}        \, d\zeta \quad \forall \phi_{k} \in \phi 
\end{gather}
If \(\mu\) is substituted with \ref{eq-aprox} the following is obtained:
\begin{gather}
\sum_{j = 1}^{N} ((\int_{\Omega} \frac{\partial^{m} \phi_{j}}{\partial \zeta^{m}} \phi_{k} \, d\zeta) c_{j}(t) - g((\int_{\Omega} \phi_k \phi_j d\zeta) c_{j}(t)))  = \int_{\Omega}  r(\zeta, t) \phi_{k}        \, d\zeta \quad \forall \phi_{k} \in \phi \label{eq-al}
\end{gather}

It is also necessary to apply the divergence theorem to the first integral term taking into account that \(y\) at \(\Gamma\) is 0.
Since \(\zeta\) is one dimensional, the divergence theorem becomes integration by parts:
\begin{gather}
\int_{\Omega} \frac{\partial^{m} \phi_{j}}{\partial \zeta^{m}} \phi_{k} \, d\zeta = - \int_{\Omega} \frac{\partial^{m-1} \phi_{j}}{\partial \zeta^{m-1}} \frac{\partial \phi_{k}}{\partial \zeta} \, d\zeta \quad \forall \phi_{k} \in \phi \label{eq-ibp}
\end{gather}
Combining \ref{eq-al} and \ref{eq-ibp} yields:
\begin{gather}
-\sum_{j = 1}^{N} ((\int_{\Omega} \frac{\partial^{m-1} \phi_{j}}{\partial \zeta^{m-1}} \frac{\partial \phi_{k}}{\partial \zeta} \, d\zeta) c_{j}(t) + g((\int_{\Omega} \phi_k \phi_j d\zeta) c_{j}(t)))  = \int_{\Omega}  r(\zeta, t) \phi_{k}        \, d\zeta \quad \forall \phi_{k} \in \phi \label{eq-fem}
\end{gather}

This formulation leads to a system of ODEs or a system of linear equations that can be solved either analytically or numerically.
\subsection{Solving the Heat Equation using FEM}
This formulation of FEM can be applied to \ref{eq-1d-h}:
\begin{gather}
\Omega = \chi \quad \Gamma = \{x_{0}, x_{n}\} \\
y(\zeta, t) = -u(x, t) \quad g(u) = -\frac{1}{\alpha} \frac{\partial u}{\partial t} \\
m = 2 \quad r(\zeta, t) = \frac{1}{\alpha} h(x,t) \\
u(x, 0) = f(x) \quad u(x_{0}, t) = 0 \quad u(x_{n}, t) = 0
\end{gather}
The set of basis functions is defined as a set of piecewise linear functions with constant step size \(\Delta x\):
\begin{gather}
    \phi_j(x)= 
\begin{cases}
    (x - x_{j-1}) / \Delta x), \quad x_{j-1} \leq x <  x_{j}\\
    (x_{j+1} - x) / \Delta x), \quad x_{j} \leq x <  x_{j + 1}\\
    0,              \quad \text{otherwise}
\end{cases}
\end{gather}
\cite{Gustafsson2011d}


The step size \(\Delta x\) is defined by \(\Delta x = \frac{x_n - x_0}{n-1}\). \label{def-delta-x}
This results in the following system of ODEs:
\begin{gather}
\sum_{j=1}^N(\int_{\chi} \phi_{j}\phi_{k}dx)\frac{dc_{j}}{dt} = \alpha \sum_{j = 1}^N(-\int_{\chi} \frac{d\phi_{j}}{dx}\frac{d\phi_{j}}{dx}dx)c_{j}(t) + \int_{\chi}h(x, t) \phi_{k} dx
\quad \forall \phi_{k} \in \phi \label{eq-heat-fem}
\end{gather}
Using matrix notation this becomes:
\begin{gather}
M^{N \times N}, K^{N \times N} \\
M\dot{c} = Kc + d \label{eq-heat-almost-ss}
\end{gather}
The matrices \(M\) and \(K\) can be easily computed (Appendix \ref{ap-mat-der}):
\begin{gather}
m_{ij} = \begin{cases}
\frac{2\Delta x}{3}, \quad k = j \\
\frac{\Delta x}{6}, \quad |k - j| = 1 \\
0, \quad otherwise 
\end{cases} \label{def-mat-a}
\quad
k_{ij} = \begin{cases}
\frac{-2\alpha}{\Delta x}, \quad k = j \\
\frac{\alpha}{\Delta x}, \quad |k - j| = 1 \\
0, otherwise
\end{cases}
\end{gather}
However it is necessary to approximate \(d\) for each point in time using numerical integration schemes.
Furthermore to solve this system of ODEs numerically an initial condition \(c_{0}\) has to be known 
\cite{Gustafsson2011b}.


It can be obtained using a least squares (LS) approach:
\begin{gather}
\sum_{j = 1}^N \langle \phi_j, \phi_k \rangle c_{j}(0) = \langle f, \phi_{k} \rangle \quad \forall \phi_k \in \phi \\
Mc_{0} = F \label{F}\\ 
c_{0} = M^{-1}F  \label{c0}
\end{gather}
\cite{Gustafsson2011c}

Observe that by multiplying \ref{eq-heat-almost-ss} with \(M^{-1}\) (\ref{ap-K-inv}) yields a system of ODEs:
\begin{gather}
\dot{c} = M^{-1}Kc + M^{-1}d \label{almost-almost-ss}
\end{gather}
This system of ODEs can be solved using an euler scheme:
\begin{gather}
c_{t+1} = \Delta t M^{-1}(Kc_{t} + d) + c_{t} \label{fem-euler}
\end{gather}
\subsubsection{Force Boundary Conidtions}
Keep in mind that vector \(d\) is time dependent and has to be recomputed for each time step.
To force the boudary condition \(u(x, t) = 0 \; x \in \Gamma\) the first and last entry of any \(c_{t}\) has to be zero:
\begin{gather}
c_{t}^{1} = 0 \quad c_{t}^{n} = 0 \quad \forall t
\end{gather}
Therefore the Matrix \(M^{-1}\) has to be adjusted:
\begin{gather}
0 = \Delta t m (Kc + d) \\
\Rightarrow m = [0, 0, 0, \dots , 0] \label{force-bound}
\end{gather}
Here \(m\) is the first row vector of \(M^{-1}\) or the last one respectively \cite{Gustafsson2011b}.
Using \ref{eq-aprox} and the computed coefficients \(c\) the function \(u(x,t)\) can be approximated.
However this is equivalent to linear interpolation between \(c_{n}\) and \(c_{n+1}\) (Appendix \ref{ap-lin-interp}).
\setlength{\parskip}{0.5\baselineskip}  % Abstand zwischen Absätzen
\rmfamily
\renewcommand*{\chapterpagestyle}{scrheadings}
\pagestyle{scrheadings}
\onehalfspacing

% ---- Literaturverzeichnis
\cleardoublepage
\renewcommand*{\chapterpagestyle}{plain}
\pagestyle{plain}
\pagenumbering{Roman}                   % Römische Seitenzahlen
\setcounter{page}{\numexpr\value{savepage}+1}
\printbibliography[title=Literaturverzeichnis]

% ---- Anhang
\appendix
\clearpage
\pagenumbering{Roman}  % römische Seitenzahlen für Anhang
\chapter{Appendix}
\section{Deriving matrices for FEM using piecewise linear functions} \label{ap-mat-der}
The so called triangle function is defined as follows:
\begin{gather}
    \phi_j(x)= 
\begin{cases}
    (x - x_{j-1}) / \Delta x), \quad x_{j-1} \leq x <  x_{j}\\
    (x_{j+1} - x) / \Delta x), \quad x_{j} \leq x <  x_{j + 1}\\
    0,              \quad \text{otherwise}
\end{cases}
\end{gather}
\cite{Gustafsson2011d}
The following integrals have to be evaluated:
\begin{gather}
\int_{\chi} \phi_{j}\phi_{k}dx \quad \forall \phi_{k} \in \phi   \label{int-1} \\
-\int_{\chi} \frac{d\phi_{j}}{dx}\frac{d\phi_{k}}{dx}dx  \quad \forall \phi_{k} \in \phi  \label{int-2}
\end{gather}
With \(\chi \subset \mathbb{R}\). 
Note that the product of two functions \(\phi_j\) and \(\phi_k\) and their derivatives is only under two conditions not zero:

1. \(k = j\)

Considering this case the integral  \ref{int-1} becomes:
\begin{gather}
\int_{x_{j-1}}^{x_{j}} \phi_j^{2} dx + \int_{x_{j}}^{x_{j+1}} \phi_j^{2} dx 
\end{gather}
Because of symmetry only one of the above integrals have to computed:
\begin{gather}
2 \int_{x_{j-1}}^{x_{j}} \phi_j^{2} dx \\
= \frac{2}{\Delta x^2} \int_{x_{j-1}}^{x_{j}} (x-x_{j-1})^{2} dx \\
\frac{2}{3 \Delta x^{2}} \left[ (x - x_{j-1})^3\right]_{x_{j-1}}^{{x_j}} = \frac{2}{3 \Delta x^{2}} \Delta x^3 = \frac{2}{3}\Delta x
\end{gather}
Integral \ref{int-2} for \(i = j\) taking symmetry into account becomes:
\begin{gather}
-\int_{x_{j-1}}^{x_{j+1}} (\frac{d \phi_j}{dx})^2 dx = -\frac{1}{\Delta x^2} \int_{x_{j-1}}^{x_{j+1}} 1 dx\\
 = -\frac{1}{\Delta x^2}  \left[ x \right]_{x_{j-1}}^{x_{j+1}} = -\frac{2}{\Delta x}
\end{gather}

2. \(|j - k| = 1\)
\ref{int-1} becomes:
\begin{gather}
\frac{1}{\Delta x^2} \int_{x_j}^{x_{j+1}} (x-x_j)(x_{j+1} - x)dx \\
= \left[ \frac{1}{2} x^2 x_{j+1} - \frac{1}{3} x^3 - x x_{j+1} x_{j} + \frac{1}{2} x^2 x_j \right]_{x_j}^{x_{x_j+1}} 
= \frac{1}{6 \Delta x^2} \Delta x^3 = \frac{1}{6} \Delta x
\end{gather}
 
Finally \ref{int-2} has to be evaluated for this condition:
\begin{gather}
-\int_{x_j}^{x_{j+1}} \frac{d\phi_{j}}{dx}\frac{d\phi_{j+1}}{dx}dx = \frac{1}{\Delta x^2} \int_{x_j}^{x_{j+1}} 1 dx =  \frac{1}{\Delta x^2} \left[ x \right]_{x_{j}}^{x_{j+1}} = \frac{1}{\Delta x}
\end{gather}
\newpage
\section{Proof that matrix A is invertible}
\label{ap-A-inv}
Let \(A_{n}\) be a matrix with \(A_{n} \in \mathbb{R}^{n \times n}\) given by:
\begin{gather}
a_{ij} = \begin{cases}
a, \quad k = j \\
b, \quad |k - j| = 1 \\
0, \quad otherwise 
\end{cases}
\end{gather}
It's determinant is given by the Laplace expansion:
\begin{gather}
det(A_n) = \sum_{j=1}^{n} (-1)^{i+j} a_{ij} M_{ij} \quad \forall i
\end{gather}
\(M_{ij}\) is the determinant of the matrix \(A'\) that is obtained by removing the \(i^{th}\) row and \(j^{th}\) column of \(A_n\).
This expression can be simplified using the definition of \(a_{ij}\):
\begin{gather}
det(A_n) = a M_{11} - b M_{12} \label{1}
\end{gather}
\(M_{11}\) is equivalent to \(det(A_{n-1})\), since the indices of rows and columns of \(A'\) are in consecutive order and \(A'\) is a \(n-1 \times n-1\) matrix :
\begin{gather}
M_{11} = det(A') \\
A' = \begin{bmatrix}
a_{22} & \dots & a_{2n}\\
\vdots & \ddots & \vdots \\
a_{n2} & \dots & a_{nn}
\end{bmatrix}
\end{gather}
\(M_{12}\) can be obtained by calculating the determinant of \(A'\) using the Laplace expansion:
\begin{gather}
A' = \begin{bmatrix}
a_{21} & a_{23} & \dots & a_{2n}\\
\vdots & \vdots & \ddots & \vdots \\
a_{n1} & a_{n3} & \dots & a_{nn}
\end{bmatrix} \\
det(A') = a_{21} \cdot det(A'') \label{an1} \\
A'' = \begin{bmatrix}
a_{33} & \dots & a_{3n}\\
\vdots & \ddots & \vdots \\
a_{n3} & \dots & a_{nn}
\end{bmatrix}
\end{gather}

In \ref{an1} only the stated term has to evaluated since all entries of the first column of the second submatrix are zero. 
Therefore the determinant is zero.
The row and column indices of \(A''\) are in consecutive order and it is a \(n-2 \times n-2\) matrix.
Therefore \(A''\) is equivalent to \(A_{n-2}\).
\ref{1} becomes:
\begin{gather}
det(A_n) = a \cdot det(A_{n-1}) - b^{2} \cdot det(A_{n-2})
\end{gather}
Furthermore this implies \(det(A_0) = 1\):
\begin{gather}
det(A_2) = a^{2} - b^{2} =  a \cdot det(A_1) - b^{2} \cdot 1 \\
\Rightarrow det(A_0) = 1
\end{gather}
Using the definition of \ref{def-mat-a} and \ref{def-delta-x} this can be seen as the following sequence:
\begin{gather}
a_0 = 1, \; a_1 = \frac{2 \Delta x}{3} \\
a_{n+1} = \frac{2 \Delta x}{3} \cdot a_{n} - \frac{\Delta x^2}{36} \cdot a_{n-1} 
\end{gather}
As described here \cite{Michael2017} a recursive sequence converges if it is monotone and has a limit.
A proof by induction shows that this sequence is monotone for \(n \geq 1\).


Base case:
\begin{gather}
a_{2} = (\frac{2 \Delta x}{3})^{2} - \frac{\Delta x^{2}}{36} = \Delta x^{2} (\frac{4}{9} - \frac{1}{36}) < \frac{2 \Delta x}{3} = a_{1}
\end{gather}
Induction step: Assuming that \(a_k < a_{k-1}\) holds, \(a_{k+1} < a_{k}\) also holds:
\begin{gather}
a_{k+1} = \frac{2 \Delta x}{3} \cdot a_{k} - \frac{\Delta x^2}{36} \cdot a_{k-1}  < \frac{2 \Delta x}{3} \cdot a_{k-1} - \frac{\Delta x^2}{36} \cdot a_{k-2} = a_{k} 
\end{gather}
The limit of this sequence is as follow:
\begin{gather}
\alpha = \lim_{n \to \infty} a_{n+1} = \lim_{n \to \infty} \Delta x \cdot \frac{2}{3} \cdot \lim_{n \to \infty} a_{n} - \lim_{n \to \infty} \Delta x^{2} 
\cdot \frac{1}{36} \cdot \lim_{n \to \infty} a_{n-1} 
= 0 \cdot \alpha - 0 \cdot \alpha = 0
\end{gather}
Since this series is monotone and converges to zero as \(n\) goes to infinity, there is no \(n \in \mathbb{N}\) for which \(a_n = 0\). 
Therefore the determinant of the matrix \(A\) defined in \ref{def-mat-a} is not zero and \(A\) is invertible.






\newpage
\end{document}
